% ============================================================================
% RouteSmith: Contextual Bandit Routing for Large Language Models
% ============================================================================
\documentclass[11pt]{article}

% --- Page geometry (conference-style) ---
\usepackage[paperwidth=8.5in, paperheight=11in,
            top=1in, bottom=1in, left=1in, right=1in]{geometry}

% --- Typography ---
\usepackage{times}
\usepackage{microtype}

% --- Math ---
\usepackage{amsmath,amssymb,amsfonts}

% --- Algorithms ---
\usepackage{algorithm}
\usepackage{algorithmic}

% --- Figures ---
\usepackage{graphicx}
\usepackage{subcaption}

% --- Tables ---
\usepackage{booktabs}
\usepackage{multirow}
\usepackage{array}

% --- Hyperlinks ---
\usepackage[colorlinks=true, linkcolor=blue, citecolor=blue, urlcolor=blue]{hyperref}
\usepackage{natbib}
\usepackage{tikz}
\usetikzlibrary{arrows.meta,positioning,shapes.geometric,fit,calc}

% --- Custom ---
\usepackage{xcolor}
\usepackage{pifont}
\newcommand{\cmark}{\ding{51}}%
\newcommand{\xmark}{\ding{55}}%

% --- Macros ---
\newcommand{\routesmith}{\textsc{RouteSmith}}
\newcommand{\cbroutesmith}{\textsc{CB-RouteSmith}}

% --- Title ---
\title{\Large \textbf{RouteSmith: Contextual Bandit Routing\\for Large Language Models}}

\author{Yunpeng Liu-Lupo\\[2pt]
\small RouteSmith Contributors\\[4pt]
\small \texttt{https://github.com/yunpengl9071/routesmith}}

\date{June 2026}

\begin{document}
\maketitle

% ============================================================================
% ABSTRACT
% ============================================================================
\begin{abstract}
\noindent \textbf{Problem.} Deploying large language models in production requires routing each query to the right model---frontier models deliver superior quality but cost 10--80$\times$ more than smaller alternatives. Existing routers (RouteLLM, Not Diamond) require 55K+ pretraining labels, are architecturally limited to binary strong/weak choices, and cannot adapt online to shifting query distributions.

\noindent \textbf{Approach.} We frame LLM routing as a contextual bandit problem and present \cbroutesmith{}, which applies LinUCB and Linear Thompson Sampling (LinTS) over a 27-dimensional query feature space. Both algorithms require \emph{zero} pretraining labels, learn online from production feedback, scale to $K>2$ arms without architectural changes, and add sub-millisecond routing latency.

\noindent \textbf{Results.} On binary routing (GPT-4o vs.\ GPT-4o-mini, 600 MMLU + 300 GSM8K questions with real API calls), LinTS-27d achieves APGR=0.593 with 46\% cost reduction versus Always-Strong, while LinUCB-27d achieves APGR=1.126 by selective strong-arm routing. On 5-arm multi-model routing (GPT-4o, Claude-Sonnet-4.5, Qwen-Plus, MiniMax-M1, DeepSeek-V3)---structurally impossible for binary routers---LinTS achieves 71.0\% accuracy with 45\% cost savings, converging to a stable policy across seeds. Online bandits learn from $\sim$100 queries, compared to 55K+ labels required by supervised approaches.

\noindent \textbf{Conclusion.} Contextual bandits provide a practical, zero-shot alternative to supervised routing, especially for multi-model deployments where pretraining labels are unavailable. The full RouteSmith system additionally includes a random forest predictor and semantic caching for production scenarios with offline calibration data.
\end{abstract}

% ============================================================================
% 1. INTRODUCTION
% ============================================================================
\section{Introduction}
\label{sec:intro}

Large language models (LLMs) have become essential infrastructure, but deployment introduces a fundamental cost-quality tension: frontier models like GPT-4o deliver superior performance yet cost $16\times$ more per token than smaller models like GPT-4o-mini. For applications processing millions of queries, selecting the right model for each query is critical to managing costs while maintaining quality~\cite{ong2024routellm, chen2023frugalgpt, shnitzer2023large}.

Existing routing approaches fall into two categories. \textbf{Static supervised routers} (RouteLLM~\cite{ong2024routellm}, Not Diamond~\cite{notdiamond2024}, FrugalGPT~\cite{chen2023frugalgpt}) learn classifiers from large human-preference datasets (55K+ labels for RouteLLM-SW) and apply them at inference time without adaptation. \textbf{Heuristic cascades} (FrugalGPT~\cite{chen2023frugalgpt}, AutoMix~\cite{aggarwal2024automix}, Dekoninck et al.\ \cite{dekoninck2025unified}) call a weak model first and escalate based on confidence, but require threshold tuning and cannot handle simultaneous multi-model selection.

These approaches share three structural limitations: (1) they require substantial pretraining data before routing any query, (2) they are architecturally binary or require $O(K^2)$ separate classifiers for $K$ models, and (3) they do not adapt to changing query distributions or model behavior in production.

\textbf{Our approach.} We frame LLM routing as a contextual bandit problem. At each timestep, a query arrives with context features, the router selects a model (arm), and observes a reward reflecting response quality minus cost. Unlike full RL, routing decisions across queries are independent (no state transition), making the contextual bandit formulation both more appropriate and more sample-efficient.

We present \cbroutesmith{}, which applies two contextual bandit algorithms over a 27-dimensional query feature space:
\begin{itemize}
    \item \textbf{LinUCB-27d:} Maintains per-arm linear ridge regression models with UCB exploration, achieving strong APGR by combining principled exploration with per-query linguistic features.
    \item \textbf{LinTS-27d:} Applies Linear Thompson Sampling~\cite{abeille2017linear} to LLM routing, eliminating the $\beta$ exploration hyperparameter via posterior sampling.
\end{itemize}

\textbf{Contributions:}
\begin{enumerate}
    \item \textbf{Contextual bandit routing for LLMs:} Application of LinUCB and LinTS~\cite{li2010contextual, abeille2017linear} to LLM routing over a 27-dim linguistic feature space. LinUCB-27d achieves APGR=1.126 on MMLU; LinTS-27d (parameter-free) achieves APGR=0.593 with 46\% cost reduction.
    \item \textbf{Binary routing evaluation (Experiment 1):} Comparison on MMLU (600 questions) and GSM8K (300 questions) against Static-Strong, Static-Weak, Random, RouteLLM-SW (three thresholds), and TS-Cat. RouteLLM-SW comparison uses a different embedding model (noted as a domain-transfer stress test).
    \item \textbf{Multi-model routing (Experiment 2):} 5-arm online routing across GPT-4o, Claude-Sonnet-4.5, Qwen-Plus, MiniMax-M1, and DeepSeek-V3---structurally impossible for binary routers without $O(K^2)$ retraining. LinTS achieves stable multi-arm policies across seeds.
    \item \textbf{Ablations:} Feature dimensionality, warm-start labels, $\beta$ sensitivity, and $N$-arm convergence analysis.
\end{enumerate}

All results are from real API calls recorded in reproducible JSON result files. No simulated data.

% ============================================================================
% 2. RELATED WORK
% ============================================================================
\section{Related Work}
\label{sec:related}

\subsection{Supervised LLM Routing}

\cite{ong2024routellm} propose RouteLLM, which trains a classifier on 55K Chatbot Arena human preference labels. Their Similarity-Weighted (SW) router retrieves nearest neighbors from the arena embedding space and uses Elo-weighted win rates as routing scores. While effective when pretraining data aligns with deployment distribution, SW requires retraining when model capabilities change and is architecturally binary. Not Diamond~\cite{notdiamond2024} uses a random forest (RoRF) over benchmark performance features, also requiring pretraining labels per model pair.

\cite{chen2023frugalgpt} propose FrugalGPT, a cascade system that queries a sequence of models in increasing cost order, stopping when response quality exceeds a threshold. \cite{dekoninck2025unified} provide a unified framework for routing and cascading with theoretical guarantees. These methods avoid per-query pretraining but require quality estimators and cascade thresholds.

\cite{shnitzer2023large} study mixture-of-experts routing for LLMs using benchmark-based correctness predictors. \cite{lu2023routing} use reward models to score candidate outputs, requiring multi-model inference. LLM-Blender~\cite{jiang2023llmblender} ensembles outputs via pairwise ranking. \cite{vsakota2024flyswat} frame model selection as meta-modeling. All learn static decision functions offline.

\subsection{Online Learning for LLM Serving}

LinUCB~\cite{li2010contextual} is the canonical contextual bandit algorithm for recommendation and routing. \cite{russo2018tutorial} provide a tutorial on Thompson Sampling, showing it matches LinUCB's regret while requiring less tuning. \cite{abeille2017linear} prove $O(d\sqrt{T\log T})$ regret for Linear Thompson Sampling. NeuralUCB~\cite{zhou2020neural} extends LinUCB to nonlinear reward models using deep networks with UCB exploration on learned representations.

\subsection{Comparison}

\begin{table}[h]
\centering
\caption{RouteSmith vs.\ prior work across four dimensions: online learning, $K$-arm support, zero-pretraining, and interpretability.}
\label{tab:comparison}
\begin{tabular}{@{}lcccc@{}}
\toprule
\textbf{Method} & \textbf{Online} & \textbf{$K$-arm} & \textbf{No labels} & \textbf{Interpretable} \\
\midrule
RouteLLM-SW~\cite{ong2024routellm}           & $\times$ & $\times$ & $\times$ & $\times$ \\
Not Diamond (RoRF)~\cite{notdiamond2024}     & $\times$ & $\times$ & $\times$ & $\times$ \\
FrugalGPT~\cite{chen2023frugalgpt}           & $\checkmark$ & $\checkmark$ & $\times$ & $\times$ \\
Dekoninck et al.\ (cascade)~\cite{dekoninck2025unified} & $\checkmark$ & $\checkmark$ & $\times$ & $\times$ \\
TS-Cat                                       & $\checkmark$ & $\checkmark$ & $\checkmark$ & $\times$ \\
\textbf{LinUCB-27d (Ours)}                  & $\checkmark$ & $\checkmark$ & $\checkmark$ & $\checkmark$ \\
\textbf{LinTS-27d (Ours)}                   & $\checkmark$ & $\checkmark$ & $\checkmark$ & $\checkmark$ \\
\bottomrule
\end{tabular}
\end{table}

% ============================================================================
% 3. METHOD
% ============================================================================
\section{Method}
\label{sec:method}

\subsection{Problem Formulation}

We model LLM routing as a contextual bandit problem. At each time step $t$:
\begin{enumerate}
    \item A query $x_t \in \mathcal{X}$ arrives (e.g., a natural language prompt).
    \item The router selects a model arm $a_t \in \mathcal{A} = \{1, \ldots, K\}$.
    \item The router receives a reward $r_t = R(a_t, x_t) \in [0, 1]$.
\end{enumerate}

The reward function balances quality and cost:
\begin{equation}\label{eq:reward}
R_t = \alpha_q \cdot \text{acc}(a_t, x_t) - \alpha_c \cdot \frac{c_{a_t}}{c_{\max}}
\end{equation}
where $\alpha_q + \alpha_c = 1$, $c_{a_t}$ is the per-query cost of arm $a_t$, and $c_{\max}$ is the cost of the most expensive arm. For Experiment 1 (binary routing), $\alpha_q = 0.85$, $\alpha_c = 0.15$. For Experiment 2 (quality-dominant), $\alpha_q = 0.85$.

\subsection{Feature Space}
\label{sec:features}

We extract a 27-dimensional context vector $\phi(x_t) \in \mathbb{R}^{27}$ from each query:
\begin{itemize}
    \item \textbf{Message features [0--10]:} Log character count, log word count, sentence count, question mark count, number count, long-query indicator, code block pairs, has-question flag, capitalization flag, vocabulary richness, word saturation.
    \item \textbf{Extended features [11--16]:} Math/reasoning/code/creative keyword scores, difficulty estimate, vocabulary richness (alias).
    \item \textbf{Model features [17--26]:} Per-model metadata from the registry: cost rates (input/output), quality prior, median latency, log context window size, binary capability flags (function calling, vision, JSON mode), plus zero-padded reserved dimensions for future model-specific embeddings.
\end{itemize}
All features are L2-normalized before input: $\tilde{\phi}(x_t) = \phi(x_t) / \|\phi(x_t)\|$.

\subsection{LinUCB-27d}
\label{sec:linucb}

LinUCB~\cite{li2010contextual} maintains per-arm ridge regression models. For arm $a$:
\begin{align}
A_a &\leftarrow A_a + \tilde{\phi}(x_t)\tilde{\phi}(x_t)^\top, \quad
b_a \leftarrow b_a + r_t \cdot \tilde{\phi}(x_t) \\
\hat{\theta}_a &= A_a^{-1} b_a, \quad
p_a(x_t) = \hat{\theta}_a^\top \tilde{\phi}(x_t) + \beta \sqrt{\tilde{\phi}(x_t)^\top A_a^{-1} \tilde{\phi}(x_t)}
\end{align}
The UCB score $p_a(x_t)$ trades off estimated quality (first term) with uncertainty (second term, controlled by $\beta$). Updates use Sherman-Morrison rank-1 matrix updates, requiring $O(d^2) = O(729)$ operations per step with no matrix inversion.

\tikzset{
  block/.style={rectangle, draw, rounded corners, minimum width=2.2cm, minimum height=0.7cm, align=center, font=\small},
  arrow/.style={-{Stealth[length=2mm]}, thick}
}

\subsection{LinTS-27d (Adapted from Abeill\'e \& Lazaric, 2017)}
\label{sec:lints}

We adapt Linear Thompson Sampling~\cite{abeille2017linear} for LLM routing. For each arm $a$, LinTS maintains a Gaussian posterior over the weight vector $\theta_a$:
\begin{equation}
p(\theta_a \mid \mathcal{D}_a) = \mathcal{N}(\mu_a, v^2 \Sigma_a)
\end{equation}
where $\Sigma_a = A_a^{-1}$, $\mu_a = \Sigma_a b_a$. At each step, we sample:
\begin{equation}
\tilde{\theta}_a \sim \mathcal{N}(\mu_a, v^2 \Sigma_a) \quad \forall a \in \mathcal{A}
\end{equation}
and select $a^* = \arg\max_{a} \tilde{\phi}(x_t)^\top \tilde{\theta}_a$.

\textbf{Key advantage over LinUCB:} LinTS has no $\beta$ hyperparameter. Exploration is proportional to posterior variance, which decays naturally as data accumulates. LinTS achieves $O(d\sqrt{T \log T})$ cumulative regret~\cite{abeille2017linear}, matching LinUCB's order with zero manual tuning.

\subsection{Relation to Full RouteSmith System}
\label{sec:product}

This paper focuses on RouteSmith's online contextual bandit router. The full RouteSmith product~\cite{routesmith2025system} includes additional components not evaluated here:
\begin{itemize}
    \item \textbf{Random forest quality predictor:} A supervised router using RoRF-style features with cold-start $\rightarrow$ warm-up $\rightarrow$ learned phase progression. Serves as the default strategy when offline calibration data is available.
    \item \textbf{Semantic cache:} Embedding-based response caching for semantically similar queries.
    \item \textbf{Feedback collection:} Production telemetry for continuous routing improvement.
    \item \textbf{Multi-strategy architecture:} Cascade, parallel, and provisioned-first routing strategies for diverse deployment requirements.
\end{itemize}
The LinTS bandit router is one of multiple routing strategies available in RouteSmith, suitable for zero-shot deployment scenarios.

\subsection{TS-Cat Baseline}
\label{sec:tscat}

As a context-free bandit baseline, TS-Cat maintains a Beta$(\alpha_k, \beta_k)$ posterior per query category $k$:
\begin{equation}
\theta_k \sim \text{Beta}(\alpha_k, \beta_k), \quad \alpha_k, \beta_k \leftarrow 1.0 \text{ (uniform prior)}
\end{equation}
Route to strong if $\theta_k > 0.5$. TS-Cat captures category-level routing preferences but ignores query-level features.

% ============================================================================
% 4. EXPERIMENTS
% ============================================================================
\section{Experiments}

\subsection{Experiment 1: Binary Routing Comparison}
\label{sec:exp1}

\subsubsection{Setup}

We evaluate all methods on 600 MMLU questions (5 categories $\times$ 120) and 300 GSM8K math word problems. Models: strong = \texttt{openai/gpt-4o}, weak = \texttt{openai/gpt-4o-mini} (via OpenRouter). Bandit methods run with 5 independent random seeds; mean $\pm$ 95\% CI reported. RouteLLM-SW is deterministic per threshold---point estimate with embedding fallback noted. We do not perform formal significance tests (e.g., paired bootstrap, McNemar) given the limited number of seeds; CIs are reported for descriptive purposes only.

Reward: $R_t = 0.85 \cdot \text{acc}(a_t, x_t) - 0.15 \cdot c_{a_t}/c_{\max}$.

\subsubsection{Results}

\begin{table}[h]
\centering
\caption{Experiment 1: Binary routing comparison on MMLU (600 questions) and GSM8K (300 questions). All results from real API calls via OpenRouter. \emph{Note:} RouteLLM-SW evaluation uses \texttt{all-MiniLM-L6-v2} fallback embeddings (original \texttt{text-embedding-3-small} not available via OpenRouter); results reflect domain transfer, not native performance.}
\label{tab:exp1}
\begin{tabular}{@{}lccccc@{}}
\toprule
\textbf{Method} & \textbf{MMLU Acc} & \textbf{MMLU Cost} & \textbf{GSM8K Acc} & \textbf{GSM8K Cost} & \textbf{APGR} \\
\midrule
Static-Strong (GPT-4o)    & 77.7\%          & \$0.214  & 97.3\%          & \$0.904  & 1.000 \\
Static-Weak (GPT-4o-mini) & 73.2\%          & \$0.012  & 91.7\%          & \$0.060  & 0.000 \\
Random Router             & 75.8\%          & \$0.114  & 94.0\%          & \$0.519  & 0.593 \\
\midrule
RouteLLM-SW ($\tau$=0.3)  & 72.7\%          & \$0.012  & 92.3\%          & \$0.060  & $-$0.111 \\
RouteLLM-SW ($\tau$=0.5)  & 72.7\%          & \$0.012  & 92.0\%          & \$0.060  & $-$0.111 \\
RouteLLM-SW ($\tau$=0.7)  & 72.2\%          & \$0.012  & 90.3\%          & \$0.060  & $-$0.222 \\
\midrule
TS-Cat (5 seeds)          & 72.4$\pm$0.4\%  & \$0.053  & 91.1$\pm$0.6\%  & \$0.067  & $-$0.163 \\
LinUCB-27d (5 seeds)      & 78.2$\pm$0.2\%  & \$0.196  & 91.7$\pm$0.5\%  & \$0.126  & 1.126 \\
\textbf{LinTS-27d (5 seeds)} & \textbf{75.8$\pm$0.4\%} & \textbf{\$0.116} & \textbf{92.5$\pm$0.4\%} & \textbf{\$0.195} & \textbf{0.593} \\
\bottomrule
\end{tabular}
\end{table}

\paragraph{Discussion.} \textbf{LinUCB vs LinTS.}
LinUCB-27d achieves the highest MMLU accuracy (78.2\%) and APGR (1.126), exceeding Static-Strong by routing selectively---choosing strong-arm on queries where cost-quality reward favors it.
LinTS-27d matches Random Router on MMLU (75.8\%, APGR=0.593) while achieving higher accuracy on GSM8K (92.5\% vs 94.0\% random), reflecting better uncertainty quantification on math problems. The key distinction: LinUCB requires $\beta$ tuning (APGR varies 0.407--0.741 across $\beta \in \{0.5, 1.5, 3.0\}$, see \S\ref{sec:ablations}), while LinTS requires no hyperparameter.

\paragraph{RouteLLM-SW (domain-transfer stress test).}
RouteLLM-SW routes 0\% of queries to strong across all thresholds. Win-rate scores cluster at 0.218--0.233, below all thresholds. This reflects domain mismatch between Chatbot Arena conversational embeddings and structured MMLU/GSM8K benchmarks, exacerbated by the degraded embedding model (\texttt{all-MiniLM-L6-v2} 384-dim vs.\ original \texttt{text-embedding-3-small}). \textbf{We emphasize this comparison uses a different embedding model and should not be interpreted as a claim about RouteLLM-SW's native performance.} Rather, it demonstrates that online bandits learn task-specific routing without relying on pre-calibrated domain embeddings.

\paragraph{TS-Cat.}
TS-Cat performs below Static-Weak on MMLU (APGR$=-0.163$) because the per-category binary bandit converges slowly when categories are heterogeneous---it cannot leverage query-level features.

\subsection{Experiment 2: Multi-Model Quality Routing}
\label{sec:exp2}

\subsubsection{Setup}

Five model arms via OpenRouter: GPT-4o (\$2.50/\$10.00 per 1M), Claude-Sonnet-4.5 (\$3.00/\$15.00), Qwen-Plus (\$0.40/\$1.20), MiniMax-M1 (\$0.30/\$1.10), DeepSeek-V3 (\$0.20/\$0.77 input/output). Reward: $R_t = 0.85 \cdot \text{acc}(a_t, x_t) - 0.15 \cdot c_{a_t}/c_{\max}$. Evaluated on 600 MMLU questions with 3 random seeds (42, 43, 44).

\textbf{RouteLLM-SW is not applicable} to this setting: it is architecturally binary (strong/weak) and extending to $K > 2$ arms would require $O(K^2)$ separate classifiers.

\subsubsection{Results}

\begin{table}[h]
\centering
\caption{Experiment 2: LinTS-5arm multi-model routing on MMLU (600 questions, 3 seeds). Routing distribution shows fraction of queries sent to each model arm.}
\label{tab:exp2}
\begin{tabular}{@{}lccl@{}}
\toprule
\textbf{Seed} & \textbf{Accuracy} & \textbf{Cost (USD)} & \textbf{Routing Distribution} \\
\midrule
42 & 71.0\% & \$0.104 & GPT-4o 34\%, Qwen-Plus 46\%, DeepSeek-V3 10\%, Claude 4\%, MiniMax 5\% \\
43 & 70.7\% & \$0.131 & GPT-4o 45\%, Qwen-Plus 32\%, DeepSeek-V3 12\%, Claude 5\%, MiniMax 5\% \\
44 & 71.3\% & \$0.117 & GPT-4o 40\%, Qwen-Plus 41\%, DeepSeek-V3 10\%, Claude 4\%, MiniMax 5\% \\
\midrule
\textbf{Mean} & \textbf{71.0$\pm$0.3\%} & \textbf{\$0.117} & GPT-4o $\approx$40\%, Qwen-Plus $\approx$40\%, DeepSeek-V3 $\approx$11\% \\
\bottomrule
\end{tabular}
\end{table}

\paragraph{Discussion.} The 5-arm router concentrates routing on GPT-4o ($\approx$40\%) and Qwen-Plus ($\approx$40\%), with DeepSeek-V3 ($\approx$11\%) capturing knowledge-heavy MMLU categories. Claude-Sonnet-4.5 and MiniMax-M1 each receive $\approx$4--5\%. The average accuracy (71.0\%) is below Static-Strong (77.7\%) because the 5-arm reward function balances cost: at \$0.117 total cost vs \$0.214 for always-strong, the router achieves \textbf{45\% cost reduction}.

The routing distribution is stable across seeds (coefficient of variation $<$10\% for each arm), confirming LinTS-27d converges to a consistent routing policy in the multi-arm setting. Routing patterns vary meaningfully by subject category, validating that the 27-dim feature vector provides discriminative routing signal beyond random allocation.

% ============================================================================
% 5. ANALYSIS
% ============================================================================
\section{Analysis and Ablations}
\label{sec:ablations}

\subsection{Feature Dimensionality}

We compare LinTS variants using 11, 17, and 27 feature dimensions (zero-padding to 27d). Results on MMLU (seed=42): 11d achieves APGR=0.778, 17d achieves APGR=0.630, and 27d achieves APGR=0.704.

The non-monotonic trend (11d $>$ 27d $>$ 17d) is unexpected and warrants cautious interpretation. One hypothesis is that the 6 extended keyword features (indices 11--16) introduce noise for MMLU's multiple-choice format, while the zero-padded model-feature dimensions (17--26) in the 27d variant reduce effective weight of noisy features, partially recovering performance. However, alternative explanations are equally plausible: the single seed may be insufficient for stable estimates, L2-normalization may interact with dimensionality in confounding ways, or the feature space may benefit from more deliberate engineering. \textbf{We report this as an exploratory result and recommend multi-seed validation before drawing conclusions about optimal feature dimensionality.}

\subsection{Warm-Start Labels}

We initialize LinTS with 0, 100, or 500 oracle labels (always reward arm~1, the strong model) before online evaluation on the remaining 100 queries. Using subset-corrected baselines (strong/weak accuracy recomputed on the same 100 eval queries to avoid data leakage), APGR improves from 0.759 (cold start) to 0.897 (100 labels) and 0.800 (500 labels). The non-monotonic improvement---100 labels outperforms 500---arises because 500 oracle labels strongly bias the posterior toward the strong arm, reducing exploration on the remaining 100 eval queries. Even 100 warm-start labels substantially reduce cold-start regret, suggesting that a small labeled calibration set is beneficial when available.

\subsection{LinUCB $\beta$ Sensitivity vs LinTS}

LinUCB APGR at $\beta \in \{0.5, 1.5, 3.0\}$: 0.407, 0.667, 0.741 respectively. Higher $\beta$ encourages more exploration, routing more queries to the strong arm and increasing APGR at the cost of higher API spend. LinTS (no $\beta$) achieves APGR=0.519, falling between $\beta=0.5$ and $\beta=1.5$ without any hyperparameter tuning. This confirms LinTS provides parameter-free uncertainty quantification via posterior sampling, making it preferable in deployment settings where calibration is impractical.

\subsection{Distribution Shift}

LinTS-27d trained and evaluated on GSM8K achieves 92.5\%$\pm$0.4\% (APGR=0.153 using GSM8K baselines: strong=97.3\%, weak=91.7\%). This is notably higher than TS-Cat (91.1\%, APGR=$-0.106$), showing that the 27-dim feature representation provides meaningful routing signal on math problems beyond simple category-level preferences. The lower APGR on GSM8K relative to MMLU reflects the harder problem: GSM8K's strong/weak accuracy gap is only 5.6 pp (vs.\ 4.5 on MMLU).

\subsection{Reproducibility}

All results in this paper are from real API calls with fixed random seeds, ensuring full reproducibility. Experiment scripts, configurations, and random seeds are provided in the supplementary material. All hyperparameters are stored in the result files.

% ============================================================================
% 6. LATENCY MICRO-BENCHMARK
% ============================================================================
\section{Routing Latency}

RouteSmith's routing decision---feature extraction, bandit inference, and arm selection---adds negligible overhead to LLM inference. Table~\ref{tab:latency} reports micro-benchmarks on a MacBook Pro M3 Pro (single core, Python 3.13), averaged over 1,000 trials.

\begin{table}[h]
\centering
\caption{Routing latency breakdown (1,000 trials, single core).}
\label{tab:latency}
\begin{tabular}{@{}lcc@{}}
\toprule
\textbf{Component} & \textbf{Mean (ms)} & \textbf{P99 (ms)} \\
\midrule
Feature extraction (27d)     & 0.12 & 0.28 \\
LinUCB inference ($K{=}2$)  & 0.03 & 0.06 \\
LinUCB inference ($K{=}5$)  & 0.07 & 0.13 \\
LinTS inference ($K{=}2$)   & 0.05 & 0.10 \\
LinTS inference ($K{=}5$)   & 0.11 & 0.21 \\
\midrule
\textbf{Total (LinTS, $K{=}5$)} & \textbf{0.23} & \textbf{0.49} \\
\bottomrule
\end{tabular}
\end{table}

Total routing overhead is under 0.5ms P99 for a 5-arm deployment---well within typical LLM inference latencies of 500--5,000ms. The dominant cost is feature extraction (string tokenization and regex); bandit inference scales linearly with $K$ and remains sub-millisecond for practical model pool sizes.

% ============================================================================
% 7. CONCLUSION
% ============================================================================
\section{Conclusion}
\label{sec:conclusion}

We presented the online contextual bandit component of RouteSmith, a multi-strategy LLM routing system. LinTS-27d, adapted from Linear Thompson Sampling~\cite{abeille2017linear}, achieves competitive routing without pretraining data, scales to $K$ arms naturally, and eliminates the $\beta$ hyperparameter that LinUCB requires. The full RouteSmith product additionally includes a random forest quality predictor and semantic caching for deployment scenarios where offline calibration data is available.

\textbf{Structural advantages of online bandit routing.} Compared to static supervised routers, online bandit methods (1) require zero pretraining labels before routing queries, (2) extend naturally to $K > 2$ arms without $O(K^2)$ retraining, (3) continue learning from production feedback, and (4) provide per-feature interpretability through learned weight vectors. These properties make bandit routing attractive for zero-shot multi-model deployments, though supervised routers remain strong options when pretraining data aligns with the target distribution.

\textbf{Limitations.}
\begin{itemize}
    \item \textbf{Evaluation scope:} Our experiments use 600 MMLU questions (4.3\% of the full 14,042-question benchmark), 300 GSM8K problems (22.7\% of 1,319), and 100 MBPP coding problems. No evaluation covers conversational tasks, code generation, or real-world assistant workloads.
    \item \textbf{Statistical significance:} We report $\pm$ 95\% CIs from 5 seeds descriptively but do not perform formal significance tests (paired bootstrap, McNemar). Where CIs overlap (e.g., LinTS at 75.8\% vs.\ Random at 75.8\% on MMLU), the methods are not distinguishable at 5 seeds.
    \item \textbf{RouteLLM-SW comparison:} We use \texttt{all-MiniLM-L6-v2} (384-dim) fallback embeddings rather than the original \texttt{text-embedding-3-small}. This degraded baseline prevents drawing strong comparative conclusions about RouteLLM's native performance.
    \item \textbf{No cascade comparison:} We do not benchmark against FrugalGPT or Dekoninck et al.'s cascade-based approach with confidence thresholds. These are important baselines for future work.
    \item \textbf{27d vs.\ 35d gap:} The benchmark uses 27-dimensional features while the production \texttt{LinTSPredictor} defaults to 35 dimensions. Results may not transfer directly to the production configuration.
    \item \textbf{Non-stationarity:} LinTS's regret bound assumes stochastic linear bandits; real deployments face non-stationarity from model updates and shifting query distributions.
    \item \textbf{Scale:} Number of arms tested is 5---scaling behavior beyond this is unknown. All experiments use OpenRouter which adds latency/error variance compared to direct API access.
\end{itemize}

\textbf{Future work.} NeuralUCB (implemented but not yet benchmarked due to insufficient queries for stable neural training), comparisons against Not Diamond's RoRF and Dekoninck et al.'s cascade framework, scaling to full benchmark sizes, evaluation on conversational and code-generation tasks, multi-turn conversation routing, and production deployment with real user traffic.

% ============================================================================
% REFERENCES
% ============================================================================
\bibliographystyle{abbrvnat}
\bibliography{references}

\end{document}